% Part: turing-machines
% Chapter: undecidability
% Section: unsolvability-decision-problem

\documentclass[../../../include/open-logic-section]{subfiles}

\begin{document}

% SEGMENT OLP-0269-S01

\olfileid{tur}{und}{dec}
\olsection{தீர்மானச் சிக்கல்}

கொடுக்கப்பட்ட ஓர் !!{sentence} ஏற்புடையதா இல்லையா என்பதைத்
தீர்மானிப்பதற்கு ஒரு விளைவுறு முறை இருந்தால், எனில் மட்டுமே
முதல்தர தருக்கம் \emph{தீர்மானிக்கத்தக்கது} என்று கூறுகிறோம்.
ஆனால் அத்தகைய முறை எதுவும் இல்லை: முதல்தர வாக்கியங்களின்
ஏற்புடைமையைத் தீர்மானிக்கும் சிக்கல் தீர்க்கவியலாதது.

இந்த முக்கியமான எதிர்மறை முடிவை நிறுவ, தீர்மானச் சிக்கலை ஒரு டியூரிங்
பொறியால் தீர்க்க முடியாது என்று நிறுவுகிறோம். அதாவது, முதல்தர
!!{sentence} ஒன்றைக் கொண்ட நாடாவில் தொடங்கப்படும் போதெல்லாம், அந்த
!!{sentence} ஏற்புடையதா இல்லையா என்பதைப் பொறுத்து, இறுதியில் நின்று
$1$ அல்லது~$0$ ஐ வெளியிடும் டியூரிங் பொறி எதுவும் இல்லை என்று
காட்டுகிறோம். சர்ச்--டியூரிங் முன்வைப்பின்படி, கணிக்கத்தக்க ஒவ்வொரு
சார்பும் டியூரிங்-கணிக்கத்தக்கது. எனவே இந்த
``ஏற்புடைமைச் சார்பு'' ஏதேனும் ஒரு வகையில் விளைவுறக்
கணிக்கத்தக்கதாக இருந்தால், அது டியூரிங்-கணிக்கத்தக்கதாக இருக்கும்.
அது டியூரிங்-கணிக்கத்தக்கதல்ல எனில், விளைவுறக் கணிக்கத்தக்கதாகவும்
இருக்க முடியாது.

தீர்மானச் சிக்கல் தீர்க்கவியலாதது என்று நிறுவுவதற்கான நமது செயலுத்தி,
நிற்றல் சிக்கலை அதற்குக் குறைப்பதாகும். இதன் பொருள் பின்வருவது:
~$e$ ஆல் விவரிக்கப்படும் டியூரிங் பொறி உள்ளீடு~$w$ இல் நின்றால்
வெளியீடு~$1$ உடனும், இல்லையெனில் வெளியீடு~$0$ உடனும் நிற்கும்
சார்பு~$h(e,w)$ டியூரிங்-கணிக்கத்தக்கதல்ல என்று நிறுவியுள்ளோம்.
முதல்தர வாக்கியங்களின் ஏற்புடைமையைத் தீர்மானிக்கும் ஒரு டியூரிங்
பொறி இருந்தால்,~$h$ ஐக் கணிக்கும் ஒரு டியூரிங் பொறியும் இருக்கும்
என்று காட்டுவோம். ஒரு டியூரிங் பொறியால்~$h$ ஐக் கணிக்க முடியாது
என்பதால், ஏற்புடைமையைத் தீர்மானிக்கும் டியூரிங் பொறியும் இருக்க முடியாது.

ஒவ்வொரு உள்ளீடு~$w$ க்கும் ஒரு டியூரிங் பொறி~$M$ க்கும்,~$M$ இன்
கட்டளைத் தொகுதியையும் உள்ளீடு~$w$ ஐயும் குறிப்பிடும்
!!a{sentence} $!T(M, w)$ ஐயும், ``$M$ இறுதியில் நிற்கிறது''
என்பதை வெளிப்படுத்தும் !!a{sentence}~$!E(M, w)$ ஐயும் பின்வருமாறு
அமையும் வகையில் விளைவுற விவரிக்க முடியும் என்று காட்டுவதே இச்செயலுத்தியின்
முதல் படி:
\begin{quote}
  $M$ உள்ளீடு~$w$ இல் நின்றால், எனில் மட்டுமே
  $\Entails !T(M, w) \lif !E(M,w)$.
\end{quote}
இவ்வாக்கியங்கள் $!T(M, w)$,~$!E(M, w)$ ஆகியவற்றை விவரிப்பதும்,
$M$ உள்ளீடு~$w$ இல் நின்றால், எனில் மட்டுமே
$!T(M, w) \lif !E(M, w)$ ஏற்புடையது என்று சரிபார்ப்பதும் நமது
நிரூபிப்பின் பெரும்பகுதியாக இருக்கும்.

\end{document}
